\documentclass[12pt]{article}

\addtolength{\textwidth}{.8in}
\addtolength{\hoffset}{-.4in}
\addtolength{\textheight}{1in}
\addtolength{\voffset}{-1in}

\usepackage{amsmath}
\usepackage{graphicx}
\usepackage{upquote} % to convert funny quotes to straight quotes
\usepackage{amsfonts}
\usepackage{paralist}
\setlength{\parindent}{15pt}

\renewcommand*{\thefootnote}{\fnsymbol{footnote}}


\begin{document}

\begin{center}
\begin{Large}
DATA 219 Notes -- Mar. 13, 2018
\end{Large}
\end{center}

\section{Accessing data from databases}

``Database" is a very loose term for ``a repository of information that can
be readily accessed." When we use that term in Data Science, we mean a data
source that is managed by a piece of software called a \textbf{DBMS}, or
\textbf{database management system}. (True, a bunch of lists, dictionaries, or
Pandas DataFrames is also ``a repository of information that can be accessed,"
but when we use the term \textbf{database} we're not talking about that.)

There are principally two reasons to learn how to access data from databases:

\begin{itemize}
\itemsep.1em
\item Sometimes, our data is simply available in that form only. We may want
to import it into Python/Pandas and deal with it using the techniques from the
first half of the semester, but in order to get it in there, we need to
(briefly) interface with the database and issue the proper commands.
\item Sometimes, our data is just flat \textit{big}. There's a limit to how
much data Spyder can store live in memory and deal with on the spot. Suppose
you had to analyze all the eBay transactions over the past two years? Or all
the tweets made in the world over the last month? That stuff can't remotely
fit in your system's memory. Hence, we have to access it through a database
which will manage the sheer size and scope for us, and work with it a little
at a time.
\end{itemize}

\subsection{Two kinds of databases}

There are essentially two kinds of databases in the world:
\textbf{relational}, and \textbf{NoSQL}.

Previously this semester, we've worked with data spread across multiple
DataFrames that we need to join together in some way. This structure is
extremely common, so much so that most of traditional database schema theory
is based directly on it. A \textbf{relational database} is a repository of
information in which the data is stored in \textbf{tables} (a table is pretty
much \textit{exactly} a DataFrame in spirit), and on which the language
\textbf{SQL} can be used to query it and manipulate it. A software system that
enables a user to create and work with relational databases is called an
\textbf{RDBMS}, or \textbf{relational database management system}.

There are many, many RDBMS's out there: Oracle, Sybase, Microsoft SQL Server,
MySQL, PostgreSQL, and literally dozens of others. They all work in very much
the same way: they store data in tables, and they are manipulated via SQL.
Ever since Ted Codd invented the relational model in 1970, this organizational
paradigm has proven to be so powerful and computationally efficient that it
has dominated the database world. Only recently have some alternate paradigms
gained some traction (so-called ``NoSQL" databases), but even so, the majority
of the world's structured data is all in relational databases of some sort.

Relational DBMS's are much more similar to each other than NoSQL systems are
to each other. The wide variety of NoSQL DBMS's (Cassandra, MongoDB, Redis,
Neo4j, \textit{etc.}) means that learning how to use one doesn't give you much
knowledge about how to use one of the others. Some store key/value pairs;
others store objects and binary data; some store information as a graph
(network). Part of the great power of the relational model is that SQL has
proven to be a flexible, broadly-applicable query language that many different
vendors have adopted and which suits many data needs.

\subsection{SQLite}

SQLite is a really slim, trim RDBMS that stores all of its data in a single
file. When you work with SQLite (or any RDBMS), though, you don't think in
terms of files, but in terms of tables. The fact that the tables are stored in
file(s) in some strange format is invisible to you as the user of the data,
and this is in fact one of the great powers of the relational model:
separating the conceptual data from the storage implementation.

\subsubsection{Interfacing to SQLite}

``Connecting" to a relational database basically means establishing a
communication path over which you can talk SQL. In Python, this is pretty
easy:

\begin{Verbatim}[fontsize=\small,samepage=true,frame=none]
import sqlite3

con = sqlite3.connect('my_local_file')
cur = con.cursor()
\end{Verbatim}

% DO I really want to mess with "cursors," though, since you can just do the
% pd.read_sql(SQL,connection) thing? (even for sqlite_master to get table
% names?)
A ``cursor" is kind of hard to get your head around, but if you think of a
``placeholder" (like the blinky cursor in Spyder or a word processing
application), you're not too far off. Imagine Google Docs, with multiple
people simultaneously editing the same document: each one has a literal cursor
showing ``where they are" in the document. Similarly, multiple programs can
access a database at the same time, and each one has a placeholder, or cursor.

What you do with a cursor is \textbf{execute an SQL command}. SQL is a new
language, but it's not nearly as complex as something like Python because it's
not a programming language, but a \textit{query language}. You don't write
programs in it; you specify queries in it. This means that the range of things
you can do in SQL is much more restricted, and hence the language is easier to
learn.

Here's the general paradigm:
\begin{Verbatim}[fontsize=\small,samepage=true,frame=none]
query_results = cur.execute("YOUR SQL STATEMENT GOES HERE")
results_list = query_results.fetchall():
\end{Verbatim}

The string inside the call to \texttt{.execute()} is called ``an embedded SQL
statement." This means we have code in one language (SQL) ``embedded" inside
another (Python). Do not be alarmed at this.

The \texttt{.fetchall()} command says ``okay, after carrying out this query on
the database, go \textit{actually} load all that data into Python." You
\textbf{don't} want to do this if the result of the query is likely to be
large. Instead, there is a \texttt{.fetchone()} command that can get one
result at a time, that you can use in a loop.

\subsubsection{Browsing the tables}

How do we know what tables exist in the database? We can run this command:

\begin{Verbatim}[fontsize=\small,samepage=true,frame=none]
query_results = cur.execute("SELECT name FROM sqlite_master WHERE type='table'")
print(query_results.fetchall())
\end{Verbatim}

to find out.

\subsubsection{Inhaling an entire table}

Once we've got the names of the tables, we often won't even need the
\texttt{cur} object any more, since Pandas provides a short cut. Suppose we
want to read an entire table into a DataFrame. \textit{If the data is small
enough to load into memory}, we can do this with a one-liner:

\begin{Verbatim}[fontsize=\small,samepage=true,frame=none]
df = pd.read_sql("select * from table_name", con)
\end{Verbatim}

This gives us a fully-fledged DataFrame, with column names, a reasonable guess
at their types, \textit{etc.} 

So one way to deal with a (smallish) SQLite database is just to connect to it
with code like the above, get all the table names, read each one into a
DataFrame, and continue in Pandas Land, forgetting all about SQLite. Now, you
can leverage all the knowledge you've learned up to this point. Bingo!

\subsubsection{(Slightly) more complex SQL}

Because the data is sometimes too large to do this, though, it's worth
learning a little SQL so you can make the RDBMS do some of the heavy lifting,
and give you trimmed results for your Python program to work with.

The workhorse is the SQL \texttt{SELECT} statement. Here are some common
variants:

\begin{itemize}
\itemsep.1em
\item To compute the \textit{number} of rows, or individual values, without
returning the entire result, use \texttt{count(*)}, \texttt{max()},
\texttt{min()}, \textit{etc.}:

\footnotesize
\begin{Verbatim}[fontsize=\small,samepage=true,frame=none]
select count(*) from table_name
select max(col1) from table_name
\end{Verbatim}
\normalsize

\item To get only the first few rows (similar to Pandas' \texttt{.head()}
operation) tack on a ``limit" expression:

\footnotesize
\begin{Verbatim}[fontsize=\small,samepage=true,frame=none]
select * from table_name limit 10
\end{Verbatim}
\normalsize

\item To get only certain columns, change the ``\texttt{*}" to a
comma-separated list of column names:

\footnotesize
\begin{Verbatim}[fontsize=\small,samepage=true,frame=none]
select col1 from table_name
select col1,col2,col3 from table_name limit 10
\end{Verbatim}
\normalsize

\item To get only certain rows, add a ``where clause" after the table name:

\footnotesize
\begin{Verbatim}[fontsize=\small,samepage=true,frame=none]
select * from table_name where col1==val1
select col2 from table_name where col1==val1 and col2>val2
select col1,col2 from table_name where col1==val1 and not col2>val2 limit 10
\end{Verbatim}
\normalsize

(The usual boolean operations work with keywords \texttt{and}, \texttt{or},
and \texttt{not}, and parentheses for grouping.)

\item The ``\texttt{in}" keyword can be useful in a ``where clause" for
querying membership in a set:

\scriptsize
\begin{Verbatim}[fontsize=\small,samepage=true,frame=none]
select * from table_name where col1 in (val1,val2)
select col1,col2 from table_name where col1==val1 and not col2 in (val1,val2,val3) limit 10
\end{Verbatim}
\normalsize

\item To get only unique rows (similar to Pandas' \texttt{.drop\_duplicates()}
functionality) use the word \texttt{distinct} before the column names:

\scriptsize
\begin{Verbatim}[fontsize=\small,samepage=true,frame=none]
select distinct * from table_name
select col1,col2 from table_name where not col1 in (val1,val2,val3) or col2!=val4 limit 5
\end{Verbatim}
\normalsize

The \texttt{distinct} option is useful in conjunction with \texttt{count()}:
\scriptsize
\begin{Verbatim}[fontsize=\small,samepage=true,frame=none]
select count(distinct col1) from table_name
\end{Verbatim}
\normalsize

\item You can have SQLite do sorting with ``\texttt{order by}":
\scriptsize
\begin{Verbatim}[fontsize=\small,samepage=true,frame=none]
select * from table_name order by col1
select distinct col1,col2 from table_name where col3==val1 order by col2,col1 desc limit 5
\end{Verbatim}
\normalsize

As you can see, you can sort on multiple columns, and sort in reverse order
with the \texttt{desc} keyword.

\item SQL supports a ``group by" operation very much like Pandas does:

\scriptsize
\begin{Verbatim}[fontsize=\small,samepage=true,frame=none]
select col1,count(*) from table_name group by col1
select col1,max(col2) from table_name where col3==val1 group by col2 order by col1 desc limit 5
\end{Verbatim}
\normalsize

\item Finally, to do a merge/join operation, use the \texttt{join} keyword:

\scriptsize
\begin{Verbatim}[fontsize=\small,samepage=true,frame=none]
select * from table_1 join table_2 on col1==col2
select col1,col2 from table_1 join table_2 on table_1.x==table_2.q where col3>val1 limit 5
\end{Verbatim}
\normalsize

The \texttt{join} keyword supports \texttt{left}, \texttt{right}, and
\texttt{outer} options.

\end{itemize}

There are lots and lots of other options as well. See the online
documentation.
\end{document}
